\chapter{Bilan général du projet}
\chapminitoc
% Ce chapitre est rare dans les rapports étudiants et c'est PRÉCISÉMENT ce
% que la grille INSA note : valeur ajoutée, facteur humain, synthèse.

\section{Synthèse des résultats au regard des objectifs}
% TODO : tableau objectifs (chap. 3) vs résultats obtenus vs statut ✓/partiel,
% avec renvoi aux chapitres.

\section{Valeur ajoutée pour ALTEN}
% TODO : plateforme de démonstration opérationnelle (démos clients/salons),
% documentation et scripts laissés (reproductibilité), briques réutilisables
% (V2V, dashboard), montée en maturité du Lab sur la 5G et l'ITS.

\section{Facteur humain et travail d'équipe}
% TODO : collaboration dans l'équipe ProtoVA, passation aux prochains
% stagiaires, communication avec les tuteurs (CDC, bilans), démos publiques.

\section{Compétences développées}
% TODO : techniques (ROS 2, réseaux 5G, normes ETSI, CAO/3D, électronique,
% web) et transverses (gestion de projet, autonomie, vulgarisation).
